%%%%%%%%%%%%%%%%%%%%%%%%%%%%%%%%%%%%%%%%%%%%%%%%%%%%%%%%%%%%%%%%%%%%%%
% Lớp 11 | Dãy số. Cấp số cộng. Cấp số nhân | Dãy số
% Dạng 1D2?1-1: Câu hỏi lý thuyết
% Số câu: 26  (sinh tự động bởi tach_cau.py; từ nay thầy sửa trực tiếp file này)
%%%%%%%%%%%%%%%%%%%%%%%%%%%%%%%%%%%%%%%%%%%%%%%%%%%%%%%%%%%%%%%%%%%%%%

%%%%%%%%%%%%%%%% TRẮC NGHIỆM 4 PHƯƠNG ÁN %%%%%%%%%%%%%%%%

% nguồn: Ok_TracNghiemLop11_Lan1
\begin{ex}%[1D2N1-1]%[ma:7033f0a3]
	Dãy số $\left(u_n\right)$ mà $u_n=(-1)^n\cdot 3$ có ba số hạng đầu tiên  lần lượt là
	 \choice
		 {$-3;-3; 3$}
		{\True $-3; 3;-3$} 
		 {$3;-3; 3$}
		{$-3;  9;-27$}
	\loigiai{ 
		Ba  số hạng đầu tiên của dãy số là $$u_1=(-1)^1\cdot 3=-3;\ u_2=(-1)^2\cdot 3=3;\ u_3=(-1)^1\cdot 3=-3.$$
	} 
\end{ex}

% nguồn: Ok_TracNghiemLop11_Lan1
\begin{ex}%[1D2N1-1]%[ma:45874c4e]
	 [1D2N1-4]
	Cho dãy  số $(a_n)$ có $a_n=\dfrac{1}{n}$. Tính chất nào sau đây của dãy số $(a_n)$ là đúng?
	 \choice
		{Hữu hạn và  tăng}
		{Hữu hạn và giảm} 
		{Vô hạn và  tăng}
		{\True  Vô hạn và giảm}
	\loigiai{ 
		Theo cách cho dãy số $(a_n)$ thì $(a_n)$ là dãy số vô hạn và khi $n$ tăng thì $\dfrac{1}{n}$ giảm.\\ 
		Vậy $(a_n)$ là dãy số vô hạn và là dãy số giảm. 
	 }
\end{ex}


%%%%%%%%%%%%%%%% ĐÚNG/SAI %%%%%%%%%%%%%%%%

% nguồn: Ok_dungsai_lop11_Lan1
\begin{ex}%[1D2H1-1]%[ma:26794089]
 Cho dãy số $(u_n)$ gồm $u_1$, $u_2$, $\ldots$, $u_n$, $\ldots$, với mọi số nguyên dương $n$.  Xét tính đúng sai các mệnh đề sau
  \choiceTF
  {Dãy số $(u_n)$ là một dãy số hữu hạn} 
  {\True Số $u_1$ gọi là số hạng  đầu}
  {Số $u_{n+1}$ được gọi là số  hạng tổng quát của dãy số}
  {\True Nếu  $u_n<u_{n+1}$ thì dãy số $(u_n)$ được gọi là dãy số tăng}
  \loigiai{
  \begin{itemchoice}
   \itemch Dãy số $(u_n)$ là  một dãy số vô hạn. Do đó mệnh đề ``Dãy số $(u_n)$ là một dãy số hữu hạn'' là mệnh đề sai.
   \itemch Số $u_1$ gọi là số hạng  đầu.
   \itemch  Số $u_n$ được gọi là số hạng tổng quát của dãy số. Do đó mệnh đề ``$u_{n+1}$ là số hạng tổng quát của dãy số'' là mệnh đề sai.
   \itemch Nếu $u_n<u_{n+1}$ thì dãy số $(u_n)$ được gọi là  dãy số tăng
  \end{itemchoice}
  }
\end{ex}

% nguồn: Ok_dungsai_lop11_Lan1
\begin{ex}%[1D2V1-1]%[ma:386542b2]
 Cho dãy số $(u_n)$ có số hạng tổng  quát $u_n=1-\dfrac{1}{n}$. Khi đó:
  \choiceTF
  {\True $u_3=\dfrac{2}{3}$} 
   {$u_7-u_8=\dfrac{1}{56}$}
   {$u_{n+1}-u_n=-\dfrac{1}{n(n+1)}$}
  {\True Dãy số  $(u_n)$ là dãy số tăng}
 \loigiai{ 
  \begin{itemchoice}
   \itemch Đúng. Vì ta có  $u_3=\dfrac{2}{3}$.
   \itemch Sai.  Vì $u_7-u_8=-\dfrac{1}{56}$.
   \itemch Sai. Vì ta có $u_{n+1}-u_n=(1-\dfrac{1}{n+1})-(1-\dfrac{1}{n})=\dfrac{1}{n}-\dfrac{1}{n-1}=\dfrac{1}{n(n+1)}>0,\,\forall n \ge  1 $.
   \itemch Đúng. Vì $u_{n+1}-u_n>0 \Rightarrow u_{n+1}>u_n,\,\forall n \in \mathbb{N}^*$. Suy  ra dãy số $(u_n)$ là dãy số tăng.
  \end{itemchoice}
 } 
\end{ex}


%%%%%%%%%%%%%%%% TỰ LUẬN %%%%%%%%%%%%%%%%

% nguồn: Ok_TuanLuan_Toan11_Lan1
\begin{bt}%[1D2Y1-1]%[ma:f09b4155]
	Chị Hương vay trả góp một khoản tiền $100$ triệu đồng và đồng ý trả dần $2$ triệu đồng mỗi tháng với lãi suất $0,8 \%$ số tiền còn lại của mỗi tháng. Gọi $A_n,$ $(n \in \mathbb{N})$ là số tiền còn nợ (triệu đồng)  của chị Hương sau $n$ tháng.
	\begin{enumerate}
		\item Tìm lần lượt  $A_0,$ $ A_1,$ $ A_2,$ $ A_3,$ $ A_4,$ $ A_5,$ $ A_6$ để tính số tiền còn nợ của chị Hương sau $6$ tháng.
		\item  Dự đoán hệ thức truy hồi đối với dãy số $\left(A_n\right)$.
	\end{enumerate}
	 \loigiai{
		\begin{enumerate}
			\item Ta có 
			\allowdisplaybreaks 
			\begin{eqnarray*}
				A_0&=&100{.}000{.}000+100{.}000{.}000 \cdot 0{,}8\%-2{.}000{.}000\\ 
				&=&100{.}000{.}000(1+0{,}8\%)-2{.}000{.}000=98{.}800{.}000.\\ 
				A_1&=&A_0+A_0 \cdot  0{,}8\%-2{.}000{.}000\\
				&=&A_0(1+0{,}8\%)-2{.}000{.}000=97{.}590{.}400. 
				\\ 
				A_2&=&A_1(1+0{,}8\%)-2{.}000{.}000=96{.}371{.}123. 
				\\ 
				 A_3&=&A_2(1+0{,}8\%)-2{.}000{.}000=95{.}142{.}092.\\
				 A_4&=&A_3(1+0{,}8\%)-2{.}000{.}000=93{.}903{.}228.\\
				 A_5&=&A_4(1+0{,}8\%)-2{.}000{.}000=92{.}654{.}454.\\
				A_6&=&A_5(1+0{,}8\%)-2{.}000{.}000=91{.}395{.}690. 
			\end{eqnarray*}
			\item Hệ thức truy hồi đối với dãy số $\left(A_n\right)$ là  $$A_n = A_{n-1}(1+0{,}8\%)-2{.}000{.}000.$$
	\end{enumerate}}
\end{bt}

% nguồn: Ok_TuanLuan_Toan11_Lan1
\begin{bt}%[1D2Y1-1]%[ma:830b0baa]
	Xét tính tăng, giảm của dãy số  $\left(u_{n}\right)$, với $u_{n}=-2 n+5$.
	 \loigiai{
		 Ta có
		$$u_{n+1}-u_{n}=[-2(n+1)+5]-(-2 n+5)=(-2  n+3)+2 n-5=-2<0,$$ tức là $u_{n+1}<u_{n},\, \forall n \in \mathbb{N}^{*}$.
	Vậy $\left(u_{n}\right)$ là  dãy số giảm.}
\end{bt}

% nguồn: Ok_TuanLuan_Toan11_Lan1
\begin{bt}%[1D2Y1-1]%[ma:0fc08dc6]
	Xét tính bị chặn của dãy số $\left(u_{n}\right)$, với $u_{n}=\dfrac{n-1}{n}$. 
	 \loigiai{
		Dãy số $\left(u_{n}\right)$ bị  chặn trên, vì $u_{n}=\dfrac{n-1}{n}=1-\dfrac{1}{n}<1,\, \forall n \in \mathbb{N}^{*}$.\\
	Dãy số $\left(u_{n}\right)$ cũng bị chặn dưới, vì $u_{n}=\dfrac{n-1}{n} \geq 0,\, \forall n \in \mathbb{N}^{*}$. Vậy dãy số $\left(u_{n}\right)$ bị chặn.} 
\end{bt}

% nguồn: Ok_TuanLuan_Toan11_Lan1
\begin{bt}%[1D2Y1-1]%[ma:afe91584]
	Xét  tính tăng, giảm của dãy số $(u_n)$, biết:
	\begin{enumEX}[a)]{3}
		\item $u_n=2 n-1$; 
		\item  $u_n=-3 n+2$;
		\item  $u_n=\dfrac{(-1)^{n-1}}{2^n}$.
	\end{enumEX}
	 \loigiai{
		\begin{enumEX}[a)]{1}
			\item Ta  có
			$$u_{n+1}-u_{n}=[2(n+1)-1]-(2 n-1)=2>0,$$ tức là $u_{n+1}>u_{n},\, \forall n \in \mathbb{N}^{*}$. Vậy $\left(u_{n}\right)$ là dãy số  tăng.
			\item  Ta có
			$$u_{n+1}-u_{n}=[-3(n+1)+2]-(-3 n+2)=-3<0,$$ tức là $u_{n+1}<u_{n},\, \forall n \in \mathbb{N}^{*}$.  Vậy $\left(u_{n}\right)$ là dãy số giảm.
			\item Bốn số hạng  đầu của dãy số là \\$u_2=\dfrac{(-1)^{1}}{2^1}=-\dfrac{1}{2};$ $u_3=\dfrac{(-1)^{2}}{2^2}=\dfrac{1}{4};$ $u_4=\dfrac{(-1)^{3}}{2^4}=-\dfrac{1}{16};$ $u_5=\dfrac{(-1)^{4}}{2^5}=\dfrac{1}{32};$ $\cdots$
			 \\
			Vậy dãy số không  tăng, không giảm.
	\end{enumEX}}
\end{bt}

% nguồn: Ok_TuanLuan_Toan11_Lan1
\begin{bt}%[1D2Y1-1]%[ma:bbe43462]
	Dãy số $(u_n)$ cho  bởi hệ thức truy hồi $u_1=1,$ $ u_n=n \cdot u_{n-1}$, với $n \geq 2$.
	\begin{enumerate}
		\item Viết năm số  hạng đầu của dãy số.
		\item Dự đoán  công thức số hạng tổng quát $u_n$.
	\end{enumerate}
	\loigiai{ 
		\begin{enumerate}
			\item Năm số hạng đầu của dãy số  là \\$u_2=2\cdot u_1=2\cdot 1=2$; $u_3=3\cdot u_2=3\cdot 2=6;$ \\$u_4=4\cdot u_3=4\cdot 6=24;$ $u_5=5\cdot u_4=5\cdot 24=120;$
			$u_6=6\cdot  u_5=6\cdot 120=720.$
			\item Số hạng tổng quát của dãy là  $u_n=n!$.
	\end{enumerate}}
\end{bt}

% nguồn: Ok_TuanLuan_Toan11_Lan1
\begin{bt}%[1D2Y1-1]%[ma:2d85fb99]
	Viết số hạng tổng quát của dãy số tăng gồm tất cả các số nguyên dương mà mỗi  số hạng của nó:
	\begin{enumEX}[a)]{2}
		\item Đều chia hết cho $3$; 
		\item Khi chia cho $4$  dư $1$.
	\end{enumEX}
	 \loigiai{
		\begin{enumEX}[a)]{2}
			\item $u_n=3k$; 
			\item  $u_n=4k+1$.
	\end{enumEX}}
\end{bt}

% nguồn: Ok_TuanLuan_Toan11_Lan1
\begin{bt}%[1D2Y1-1]%[ma:a37e5614]
	\immini{Cho hình tứ diện $ABCD$. Trong các mặt phẳng chứa các mặt của hình tứ diện, hãy cho  biết:
	\begin{listEX}[1]
		\item Đường thẳng $A B$ cắt các mặt  phẳng nào;
		\item Đường thẳng $A B$ nằm trong các  mặt phẳng nào.
	\end{listEX}}{\begin{tikzpicture}[scale=.8,>=stealth,font=\footnotesize,line join=round,line cap=round]
		\tkzDefPoints{0/0/B,1.3/-1.6/C,4.5/0/D,1/3.5/A} 
		\tkzDrawPolygon  (A,B,C,D)
		\tkzDrawSegments  (A,C)
		\tkzDrawSegments[dashed](B,D) 
		 \tkzDrawPoints[fill=black,size=1.5](A,B,C,D)
		\tkzLabelPoints[above](A) 
		\tkzLabelPoints[below](C) 
		\tkzLabelPoints[left](B) 
		 \tkzLabelPoints[right](D)
	\end{tikzpicture}
	} 
	 \loigiai{
		\begin{listEX}[1]
			\item Đường thẳng $AB$ cắt các mặt phẳng $(ACD)$  và $(BCD)$.
			\item Đường thẳng $AB$ nằm trong các mặt phẳng $(ABC)$ và  $(ABD)$.
		\end{listEX}
	} 
\end{bt}

% nguồn: Ok_TuanLuan_Toan11_Lan1
\begin{bt}%[1D2Y1-1]%[ma:f68920c6]
	Cho  hai hình bình hành $ABCD$ và $ABEF$ không cùng nằm trong một mặt phẳng.
	\begin{enumerate}
		\item Chứng minh  rằng $(AFD)\parallel(BEC)$.
		\item Gọi $M$ là trọng tâm của tam giác $ABE$. Gọi $(P)$ là mặt phẳng đi qua $M$  và song song với mặt phẳng $(AFD)$. Lấy $N$ là giao điểm của $(P)$ và $A C$. Tính $\dfrac{AN}{NC}$.
	\end{enumerate}
	\loigiai{ 
	\immini{\begin{listEX}
		\item Ta có $AD\parallel BC$, suy  ra $AD\parallel (BCE)$. Tương tự $AF\parallel (BCE)$.\\
		 Khi đó $(ADF)\parallel (BCE)$.
		 \item Ta có $(P)\parallel(AFD)$ nên $(P)\cap (ABF)=MH$ với $M\in AB$ và $MH\parallel AF$.\\
		Vì $M$ là trọng tâm của tam  giác $ABE$ nên $\dfrac{BM}{BO}=\dfrac{2}{3}$ với $O$ là tâm hình bình hành $ABEF$.\\
		Vì $ABEF$ là hình bình hành nên $\dfrac{BO}{BF}=\dfrac{1}{2}$. Suy ra $\dfrac{BM}{BD}=\dfrac{1}{3}$. Do đó $\dfrac{BM}{MF}=\dfrac{1}{2}$.\\ 
		Khi đó $\dfrac{BH}{AH}=\dfrac{1}{2}$ (vì  $AF\parallel MH$).
	\end{listEX}}{
	\begin{tikzpicture}[scale=0.8,font=\footnotesize,line join=round,line cap=round,>=stealth]
		\coordinate (C) at  (2,-1);
		 \coordinate (F) at (-1,4);
		\coordinate (B) at (4,1); 
		\coordinate (A)  at (-2,1);
		 \coordinate (D) at (-4,-1);
		\coordinate (E) at  (5,4);
		\coordinate (O) at ($(A)!0.5!(E)$); 
		\coordinate (M)  at (barycentric cs:A=1,B=1,E=1);
		\coordinate (H)  at ($(B)!1/3!(A)$);
		\coordinate (N) at  ($(A)!2/3!(C)$);
		\draw (D)--(C)--(B)--(E)--(F)--(D); 
		\draw  (C)--(E);
		 \draw[dashed] (A)--(F)
			(A)--(B) 
			(A)--(D); 
		 \draw[dashed] (F)--(B);
		\draw[dashed]  (A)--(C);
		 \draw[dashed] (M)--(H)
			 (H)--(N);
		\foreach \p in {A,B,C,D,E,F,M,H,N,O} 
		\fill (\p) circle  (1pt);
		\node[above  right] at (H){$H$};
		 \node[below] at (N){$N$};
		\node[above] at (O){$O$}; 
		\node[above]  at (F){$F$};
		\node[left] at  (A){$A$};
		\node[right]  at (B){$B$};
		\node[above]  at (E){$E$};
		\node[above]  at (M){$M$};
		 \node[right] at (C){$C$};
		 \node[left] at (D){$D$};
	\end{tikzpicture}
	} 
	\noindent $(P)\parallel(AFD)$ nên $(P)\cap (ABC)=NH$ với $N\in AC$ và $NH\parallel  AD$.\\
Do  đó $\dfrac{AN}{NC}=\dfrac{AH}{BH}=2$.}
\end{bt}

% nguồn: Ok_TuanLuan_Toan11_Lan1
\begin{bt}%[1D2Y1-1]%[ma:5cf3f8cf]
	Ông An gửi tiết kiệm $100$ triệu đồng kì hạn $1$ tháng với lãi suất $6 \%$ một năm theo hình thức tính lãi kép.  Số tiền (triệu đồng) của ông An thu được sau $n$ tháng được cho bởi công thức
	$$ 
	 A_n=100\left(1+\dfrac{0{,}06}{12}\right)^n .
	$$ 
	\begin{enumerate}
		\item Tìm số tiền ông An nhận được sau  tháng thứ nhất, sau tháng thứ hai.
		\item Tìm  số tiền ông An nhận được sau 1 năm.
	\end{enumerate}
	\loigiai{ 
		\begin{enumerate}
			\item Số tiền ông An nhận được sau tháng thứ nhất là $ 
			A_n=100\left(1+\dfrac{0{,}06}{12}\right)^1 =100{,}5$ (triệu đồng), sau tháng thứ  hai là $
			 A_n=100\left(1+\dfrac{0{,}06}{12}\right)^2
			=101{,}0025$ (triệu  đồng).
			\item Số tiền ông An nhận được  sau 1 năm là $
			 A_n=100\left(1+\dfrac{0{,}06}{12}\right)^{12} =106{,}1678$ (triệu đồng).
	\end{enumerate}}
\end{bt}

% nguồn: Ok_TuanLuan_Toan11_Lan1
\begin{bt}%[1D2Y1-1]%[ma:bf634069]
	Trong các dãy số $(u_n)$ sau, dãy số nào bị chặn  dưới, bị chặn trên, bị chặn?
	\begin{enumEX}[a)]{4}
		\item  $u_n=n-1$;
		 \item $u_n=\dfrac{n+1}{n+2}$;
		\item $u_n=\sin  n$;
		 \item $u_n=(-1)^{n-1} n^2$.
	\end{enumEX}
	\loigiai{\begin{enumEX}[a)]{1}
		\item  Dãy số $u_n=n-1$ bị chặn dưới, vì $u_{n}\ge 0,\, \forall n \in \mathbb{N}^{*}$.
		\item Dãy số $u_n=\dfrac{n+1}{n+2}$  bị chặn trên, vì $u_n=1-\dfrac{1}{n+2}<1,\, \forall n \in \mathbb{N}^{*}$.\\
		Dãy số $\left(u_{n}\right)$ cũng bị chặn dưới, vì $u_n=\dfrac{n+1}{n+2} \geq \dfrac{2}{3},\, \forall  n \in \mathbb{N}^{*}$.\\
		 Vậy dãy số $\left(u_{n}\right)$ bị chặn.
		\item Dãy số $u_n=\sin n$ bị chặn trên, bị chặn dưới, vì $-1\le \sin n  \le1,\, \forall n \in \mathbb{N}^{*}$.\\
		 Vậy dãy số $\left(u_{n}\right)$ bị chặn.
		\item Bốn  số hạng đầu của dãy số là \\$u_2=(-1)^1\cdot 2^2=-4;$ $u_3=(-1)^2\cdot 3^2=9;$ $u_4=(-1)^3 \cdot 4^2=-16;$ $u_5=(-1)^4\cdot 5^2=25;$ $\cdots$
		 \\
		Vậy dãy số không  bị chặn.
\end{enumEX}}
\end{bt}

% nguồn: Ok_TuanLuan_Toan11_Lan1
\begin{bt}%[1D2Y1-1]%[ma:7475a090]
	Viết năm số hạng đầu và số hạng thứ 100 của các dãy số $(u_n)$ có số  hạng tổng quát cho bởi:
	\begin{enumEX}[a)]{3}
		\item  $u_n=3 n-2$;
		\item $u_n=3 \cdot 2^n$; 
		\item  $u_n=\left(1+\dfrac{1}{n}\right)^{n}$.
	\end{enumEX}
	 \loigiai{
		\begin{enumEX}[a)]{1}
			\item $u_1=3\cdot 1-2=1;$ $u_2=3\cdot 2-2=4;$ $u_3=3\cdot 3-2=7;$  $u_4=3\cdot 4-2=10;$ $u_5=3\cdot 5-2=13;$ $u_{100}=3\cdot 100-2=298.$
			\item $u_1=3 \cdot 2^1=6;$ $u_2=3 \cdot 2^2=12;$ $u_3=3 \cdot 2^3=24;$ $u_4=3 \cdot 2^4=48;$ $u_5=3 \cdot 2^5=96;$  $u_{100}=3 \cdot 2^{100}.$
			\item $u_1=\left(1+\dfrac{1}{1}\right)^{1}=2$; $u_2=\left(1+\dfrac{1}{2}\right)^{2}=\dfrac{9}{4}$; $u_3=\left(1+\dfrac{1}{3}\right)^{3}=\dfrac{64}{27}$; $u_4=\left(1+\dfrac{1}{4}\right)^{4}=\dfrac{625}{256}$; $u_5=\left(1+\dfrac{1}{5}\right)^{5}=\dfrac{1296}{625}$;  $u_{100}=\left(1+\dfrac{1}{100}\right)^{100}$.
	\end{enumEX}}
\end{bt}

% nguồn: Ok_TuanLuan_Toan11_Lan2
\begin{ex}%[1D2Y1-1]%[ma:a066a66c]
	Chị Hương vay trả góp một khoản tiền $100$ triệu đồng và đồng ý trả dần $2$ triệu đồng mỗi tháng với lãi suất $0,8 \%$ số tiền còn lại của mỗi tháng. Gọi $A_n,$ $(n \in \mathbb{N})$ là số tiền còn nợ (triệu đồng) của chị Hương  sau $n$ tháng.
	\begin{enumerate}
		\item Tìm lần lượt $A_0,$ $  A_1,$ $ A_2,$ $ A_3,$ $ A_4,$ $ A_5,$ $ A_6$ để tính số tiền còn nợ của chị Hương sau $6$ tháng.
		\item Dự đoán hệ thức truy hồi đối với dãy số $\left(A_n\right)$. 
	\end{enumerate}
	 \loigiai{
		\begin{enumerate}
			\item  Ta có
			\allowdisplaybreaks 
			\begin{eqnarray*}
				A_0&=&100{.}000{.}000+100{.}000{.}000  \cdot 0{,}8\%-2{.}000{.}000\\
				 &=&100{.}000{.}000(1+0{,}8\%)-2{.}000{.}000=98{.}800{.}000.\\
				A_1&=&A_0+A_0  \cdot 0{,}8\%-2{.}000{.}000\\
				&=&A_0(1+0{,}8\%)-2{.}000{.}000=97{.}590{.}400. 
				 \\
				A_2&=&A_1(1+0{,}8\%)-2{.}000{.}000=96{.}371{.}123. 
				\\ 
				A_3&=&A_2(1+0{,}8\%)-2{.}000{.}000=95{.}142{.}092.\\ 
				 A_4&=&A_3(1+0{,}8\%)-2{.}000{.}000=93{.}903{.}228.\\
				 A_5&=&A_4(1+0{,}8\%)-2{.}000{.}000=92{.}654{.}454.\\
				 A_6&=&A_5(1+0{,}8\%)-2{.}000{.}000=91{.}395{.}690.
			\end{eqnarray*}
			\item Hệ thức truy hồi đối với dãy  số $\left(A_n\right)$ là $$A_n = A_{n-1}(1+0{,}8\%)-2{.}000{.}000.$$
	\end{enumerate}}
\end{ex}

% nguồn: Ok_TuanLuan_Toan11_Lan2
\begin{ex}%[1D2Y1-1]%[ma:0d7b3d9d]
	Xét tính tăng, giảm  của dãy số $\left(u_{n}\right)$, với $u_{n}=-2 n+5$.
	 \loigiai{
		Ta  có
		$$u_{n+1}-u_{n}=[-2(n+1)+5]-(-2 n+5)=(-2 n+3)+2 n-5=-2<0,$$ tức  là $u_{n+1}<u_{n},\, \forall n \in \mathbb{N}^{*}$.
	Vậy $\left(u_{n}\right)$ là  dãy số giảm.}
\end{ex}

% nguồn: Ok_TuanLuan_Toan11_Lan2
\begin{ex}%[1D2Y1-1]%[ma:59410751]
	 Xét tính bị chặn của dãy số $\left(u_{n}\right)$, với $u_{n}=\dfrac{n-1}{n}$.
	\loigiai{ 
		Dãy số $\left(u_{n}\right)$ bị  chặn trên, vì $u_{n}=\dfrac{n-1}{n}=1-\dfrac{1}{n}<1,\, \forall n \in \mathbb{N}^{*}$.\\
	Dãy số $\left(u_{n}\right)$ cũng bị chặn dưới, vì $u_{n}=\dfrac{n-1}{n}  \geq 0,\, \forall n \in \mathbb{N}^{*}$. Vậy dãy số $\left(u_{n}\right)$ bị chặn.}
\end{ex}

% nguồn: Ok_TuanLuan_Toan11_Lan2
\begin{ex}%[1D2Y1-1]%[ma:8a94a6a8]
	Xét tính tăng, giảm của dãy số  $(u_n)$, biết:
	\begin{enumEX}[a)]{3}
		\item  $u_n=2 n-1$;
		\item $u_n=-3 n+2$; 
		\item $u_n=\dfrac{(-1)^{n-1}}{2^n}$. 
	\end{enumEX}
	 \loigiai{
		\begin{enumEX}[a)]{1}
			 \item Ta có
			$$u_{n+1}-u_{n}=[2(n+1)-1]-(2 n-1)=2>0,$$ tức là  $u_{n+1}>u_{n},\, \forall n \in \mathbb{N}^{*}$. Vậy $\left(u_{n}\right)$ là dãy số tăng.
			\item Ta có 
			$$u_{n+1}-u_{n}=[-3(n+1)+2]-(-3 n+2)=-3<0,$$ tức là $u_{n+1}<u_{n},\, \forall n \in \mathbb{N}^{*}$. Vậy $\left(u_{n}\right)$ là dãy số  giảm.
			\item Bốn số hạng đầu của dãy số là  \\$u_2=\dfrac{(-1)^{1}}{2^1}=-\dfrac{1}{2};$ $u_3=\dfrac{(-1)^{2}}{2^2}=\dfrac{1}{4};$ $u_4=\dfrac{(-1)^{3}}{2^4}=-\dfrac{1}{16};$ $u_5=\dfrac{(-1)^{4}}{2^5}=\dfrac{1}{32};$ $\cdots$
			 \\
			Vậy dãy số không tăng, không giảm. 
	\end{enumEX}}
\end{ex}

% nguồn: Ok_TuanLuan_Toan11_Lan2
\begin{ex}%[1D2Y1-1]%[ma:71625e49]
	Dãy số $(u_n)$ cho bởi hệ thức truy hồi $u_1=1,$ $ u_n=n \cdot u_{n-1}$,  với $n \geq 2$.
	\begin{enumerate}
		\item Viết năm  số hạng đầu của dãy số.
		\item  Dự đoán công thức số hạng tổng quát $u_n$.
	\end{enumerate}
	 \loigiai{
		\begin{enumerate}
			\item Năm  số hạng đầu của dãy số là \\$u_2=2\cdot u_1=2\cdot 1=2$; $u_3=3\cdot u_2=3\cdot 2=6;$ \\$u_4=4\cdot u_3=4\cdot 6=24;$ $u_5=5\cdot u_4=5\cdot 24=120;$
			$u_6=6\cdot  u_5=6\cdot 120=720.$
			\item Số hạng tổng quát của dãy là $u_n=n!$. 
	\end{enumerate}}
\end{ex}

% nguồn: Ok_TuanLuan_Toan11_Lan2
\begin{ex}%[1D2Y1-1]%[ma:439d7877]
	Viết số hạng tổng  quát của dãy số tăng gồm tất cả các số nguyên dương mà mỗi số hạng của nó:
	\begin{enumEX}[a)]{2}
		\item Đều chia hết cho  $3$;
		\item Khi chia cho $4$  dư $1$.
	\end{enumEX}
	 \loigiai{
		\begin{enumEX}[a)]{2}
			\item $u_n=3k$; 
			 \item $u_n=4k+1$.
	\end{enumEX}}
\end{ex}

% nguồn: Ok_TuanLuan_Toan11_Lan2
\begin{ex}%[1D2Y1-1]%[ma:3aeccea1]
	\immini{Cho hình tứ diện $ABCD$. Trong các mặt phẳng chứa  các mặt của hình tứ diện, hãy cho biết:
	\begin{listEX}[1]
		\item Đường thẳng $A B$ cắt các  mặt phẳng nào;
		\item Đường thẳng $A B$ nằm trong các mặt phẳng  nào.
	\end{listEX}}{\begin{tikzpicture}[scale=.8,>=stealth,font=\footnotesize,line join=round,line cap=round]
		\tkzDefPoints{0/0/B,1.3/-1.6/C,4.5/0/D,1/3.5/A} 
		\tkzDrawPolygon  (A,B,C,D)
		\tkzDrawSegments  (A,C)
		\tkzDrawSegments[dashed](B,D) 
		 \tkzDrawPoints[fill=black,size=1.5](A,B,C,D)
		 \tkzLabelPoints[above](A)
		 \tkzLabelPoints[below](C)
		\tkzLabelPoints[left](B) 
		\tkzLabelPoints[right](D) 
	\end{tikzpicture}
	 }
	\loigiai{ 
		\begin{listEX}[1]
			\item Đường thẳng $AB$  cắt các mặt phẳng $(ACD)$ và $(BCD)$.
			\item  Đường thẳng $AB$ nằm trong các mặt phẳng $(ABC)$ và $(ABD)$.
		\end{listEX}
	} 
\end{ex}

% nguồn: Ok_TuanLuan_Toan11_Lan2
\begin{ex}%[1D2Y1-1]%[ma:5a5ac4fc]
	Cho hai hình bình hành $ABCD$  và $ABEF$ không cùng nằm trong một mặt phẳng.
	\begin{enumerate}
		\item Chứng minh  rằng $(AFD)\parallel(BEC)$.
		\item Gọi $M$ là  trọng tâm của tam giác $ABE$. Gọi $(P)$ là mặt phẳng đi qua $M$ và song song với mặt phẳng $(AFD)$. Lấy $N$ là giao điểm của $(P)$ và $A C$. Tính $\dfrac{AN}{NC}$.
	\end{enumerate}
	 \loigiai{
	\immini{\begin{listEX}
		\item Ta có $AD\parallel BC$, suy ra $AD\parallel (BCE)$. Tương tự  $AF\parallel (BCE)$.\\
		Khi đó  $(ADF)\parallel (BCE)$.
		 \item Ta có $(P)\parallel(AFD)$ nên $(P)\cap (ABF)=MH$ với $M\in AB$ và $MH\parallel AF$.\\
		Vì $M$ là trọng tâm của tam giác $ABE$ nên $\dfrac{BM}{BO}=\dfrac{2}{3}$ với $O$ là tâm hình  bình hành $ABEF$.\\
		 Vì $ABEF$ là hình bình hành nên $\dfrac{BO}{BF}=\dfrac{1}{2}$. Suy ra $\dfrac{BM}{BD}=\dfrac{1}{3}$. Do đó $\dfrac{BM}{MF}=\dfrac{1}{2}$.\\
		Khi đó  $\dfrac{BH}{AH}=\dfrac{1}{2}$ (vì $AF\parallel MH$).
	\end{listEX}}{
	\begin{tikzpicture}[scale=0.8,font=\footnotesize,line join=round,line cap=round,>=stealth]
		\coordinate  (C) at (2,-1);
		 \coordinate (F) at (-1,4);
		\coordinate  (B) at (4,1);
		\coordinate  (A) at (-2,1);
		\coordinate  (D) at (-4,-1);
		 \coordinate (E) at (5,4);
		\coordinate (O) at  ($(A)!0.5!(E)$);
		\coordinate (M) at (barycentric  cs:A=1,B=1,E=1);
		\coordinate  (H) at ($(B)!1/3!(A)$);
		\coordinate  (N) at ($(A)!2/3!(C)$);
		 \draw (D)--(C)--(B)--(E)--(F)--(D);
		\draw (C)--(E); 
		 \draw[dashed] (A)--(F)
			(A)--(B) 
			(A)--(D); 
		\draw[dashed]  (F)--(B);
		\draw[dashed]  (A)--(C);
		\draw[dashed]  (M)--(H)
			 (H)--(N);
		\foreach \p in {A,B,C,D,E,F,M,H,N,O} 
		\fill  (\p) circle (1pt);
		\node[above right] at  (H){$H$};
		\node[below]  at (N){$N$};
		 \node[above] at (O){$O$};
		\node[above]  at (F){$F$};
		 \node[left] at (A){$A$};
		\node[right] at  (B){$B$};
		\node[above] at  (E){$E$};
		\node[above] at (M){$M$}; 
		\node[right]  at (C){$C$};
		\node[left]  at (D){$D$};
	\end{tikzpicture}
	} 
	\noindent $(P)\parallel(AFD)$ nên $(P)\cap (ABC)=NH$ với $N\in AC$ và $NH\parallel AD$.\\ 
Do đó $\dfrac{AN}{NC}=\dfrac{AH}{BH}=2$.} 
\end{ex}

% nguồn: Ok_TuanLuan_Toan11_Lan2
\begin{ex}%[1D2Y1-1]%[ma:4dbb9c8f]
	Ông An gửi tiết kiệm $100$ triệu đồng kì hạn $1$ tháng với lãi suất $6 \%$ một năm theo hình thức tính lãi kép. Số tiền (triệu đồng) của ông An thu được sau $n$ tháng được cho bởi công thức 
	$$ 
	 A_n=100\left(1+\dfrac{0{,}06}{12}\right)^n .
	 $$
	\begin{enumerate}
		 \item Tìm số tiền ông An nhận được sau tháng thứ nhất, sau tháng thứ hai.
		 \item Tìm số tiền ông An nhận được sau 1 năm.
	\end{enumerate}
	\loigiai{ 
		\begin{enumerate}
			\item Số tiền ông An nhận được sau tháng thứ nhất là $ 
			A_n=100\left(1+\dfrac{0{,}06}{12}\right)^1 =100{,}5$ (triệu đồng), sau tháng thứ hai là $ 
			A_n=100\left(1+\dfrac{0{,}06}{12}\right)^2 
			 =101{,}0025$ (triệu đồng).
			\item Số tiền ông An nhận được sau 1 năm là $ 
			A_n=100\left(1+\dfrac{0{,}06}{12}\right)^{12} =106{,}1678$  (triệu đồng).
	\end{enumerate}}
\end{ex}

% nguồn: Ok_TuanLuan_Toan11_Lan2
\begin{ex}%[1D2Y1-1]%[ma:a3b35e9b]
	Trong các dãy số  $(u_n)$ sau, dãy số nào bị chặn dưới, bị chặn trên, bị chặn?
	\begin{enumEX}[a)]{4}
		\item  $u_n=n-1$;
		\item  $u_n=\dfrac{n+1}{n+2}$;
		\item $u_n=\sin  n$;
		\item $u_n=(-1)^{n-1} n^2$. 
	\end{enumEX}
	\loigiai{\begin{enumEX}[a)]{1}
		\item Dãy số $u_n=n-1$ bị chặn dưới, vì $u_{n}\ge 0,\, \forall n  \in \mathbb{N}^{*}$.
		 \item Dãy số $u_n=\dfrac{n+1}{n+2}$ bị chặn trên, vì $u_n=1-\dfrac{1}{n+2}<1,\, \forall n \in \mathbb{N}^{*}$.\\
		 Dãy số $\left(u_{n}\right)$ cũng bị chặn dưới, vì $u_n=\dfrac{n+1}{n+2} \geq \dfrac{2}{3},\, \forall n \in \mathbb{N}^{*}$.\\
		Vậy dãy số $\left(u_{n}\right)$  bị chặn.
		\item Dãy số $u_n=\sin  n$ bị chặn trên, bị chặn dưới, vì $-1\le \sin n \le1,\, \forall n \in \mathbb{N}^{*}$.\\
		Vậy dãy  số $\left(u_{n}\right)$ bị chặn.
		\item Bốn số hạng  đầu của dãy số là \\$u_2=(-1)^1\cdot 2^2=-4;$ $u_3=(-1)^2\cdot 3^2=9;$ $u_4=(-1)^3 \cdot 4^2=-16;$ $u_5=(-1)^4\cdot 5^2=25;$ $\cdots$
		\\ 
		Vậy dãy số  không bị chặn.
\end{enumEX}}
\end{ex}

% nguồn: Ok_TuanLuan_Toan11_Lan2
\begin{ex}%[1D2Y1-1]%[ma:65e366d0]
	Viết năm số hạng đầu  và số hạng thứ 100 của các dãy số $(u_n)$ có số hạng tổng quát cho bởi:
	\begin{enumEX}[a)]{3}
		\item $u_n=3 n-2$; 
		\item $u_n=3  \cdot 2^n$;
		\item $u_n=\left(1+\dfrac{1}{n}\right)^{n}$. 
	\end{enumEX}
	 \loigiai{
		\begin{enumEX}[a)]{1}
			 \item $u_1=3\cdot 1-2=1;$ $u_2=3\cdot 2-2=4;$ $u_3=3\cdot 3-2=7;$ $u_4=3\cdot 4-2=10;$ $u_5=3\cdot 5-2=13;$ $u_{100}=3\cdot 100-2=298.$
			\item $u_1=3 \cdot 2^1=6;$ $u_2=3 \cdot 2^2=12;$ $u_3=3 \cdot  2^3=24;$ $u_4=3 \cdot 2^4=48;$ $u_5=3 \cdot 2^5=96;$ $u_{100}=3 \cdot 2^{100}.$
			\item $u_1=\left(1+\dfrac{1}{1}\right)^{1}=2$; $u_2=\left(1+\dfrac{1}{2}\right)^{2}=\dfrac{9}{4}$; $u_3=\left(1+\dfrac{1}{3}\right)^{3}=\dfrac{64}{27}$;  $u_4=\left(1+\dfrac{1}{4}\right)^{4}=\dfrac{625}{256}$; $u_5=\left(1+\dfrac{1}{5}\right)^{5}=\dfrac{1296}{625}$; $u_{100}=\left(1+\dfrac{1}{100}\right)^{100}$.
	\end{enumEX}}
\end{ex}
